\section*{Sets and Indices}

\begin{itemize}
    \item $L = \{1,\dots,n\}$: set of locations, with depot at location $1$.
    \item $C = L \setminus \{1\}$: set of customer locations.
    \item $K = \{1,2\}$: set of candidate service days.
    \item Indices:
    \begin{itemize}
        \item $i,j \in L$ for locations,
        \item $d \in K$ for days.
    \end{itemize}
\end{itemize}

\section*{Parameters}

\begin{itemize}
    \item $n$ (code: \texttt{num\_customers}): number of locations including the depot; here $n=7$.
    \item $D_{ij}$ (code: \texttt{D[i][j]}): travel distance from location $i$ to location $j$, with $D_{ii}=0$ and symmetric values taken from the distance table.
    \item $M$ (code: \texttt{max\_route\_length\_day}): maximum allowed route length for each day.
    \item $f$ (code: \texttt{route\_cost\_day}): fixed cost incurred if day $d$ is operated.
\end{itemize}

\section*{Decision Variables}

\begin{itemize}
    \item $x_{ijd}$ (code: \texttt{x[i,j,d]}) $\in \{0,1\}$ for $i,j \in L$, $i \neq j$, $d \in K$:
    equals $1$ if arc $(i,j)$ is used on day $d$.
    \item $y_{id}$ (code: \texttt{y[i,d]}) $\in \{0,1\}$ for $i \in C$, $d \in K$:
    equals $1$ if customer $i$ is assigned to day $d$.
    \item $z_d$ (code: \texttt{z[d]}) $\in \{0,1\}$ for $d \in K$:
    equals $1$ if day $d$ is active.
    \item $u_{id}$ (code: \texttt{u[i,d]}) for $i \in C$, $d \in K$:
    MTZ ordering variable with
    \[
    1 \le u_{id} \le n-1.
    \]
\end{itemize}

\section*{Objective}

Minimize total travel distance plus the fixed cost of each active day:
\[
\min \;
\sum_{d \in K} \sum_{\substack{i \in L \\ }} \sum_{\substack{j \in L \\ j \neq i}} D_{ij}\, x_{ijd}
\;+\;
\sum_{d \in K} f\, z_d .
\]

\section*{Constraints}

Each customer is assigned to exactly one day:
\[
\sum_{d \in K} y_{id} = 1
\qquad \forall i \in C.
\]

If a customer is assigned to a day, exactly one outgoing arc and one incoming arc are used for that customer on that day:
\[
\sum_{\substack{j \in L \\ j \neq i}} x_{ijd} = y_{id}
\qquad \forall i \in C,\; d \in K,
\]
\[
\sum_{\substack{j \in L \\ j \neq i}} x_{jid} = y_{id}
\qquad \forall i \in C,\; d \in K.
\]

For each day, the depot has exactly one departure and one return if the day is active, and none otherwise:
\[
\sum_{j \in C} x_{1jd} = z_d
\qquad \forall d \in K,
\]
\[
\sum_{j \in C} x_{j1d} = z_d
\qquad \forall d \in K.
\]

A customer can only be assigned to an active day:
\[
y_{id} \le z_d
\qquad \forall i \in C,\; d \in K.
\]

MTZ subtour elimination within each day, over customer nodes only:
\[
u_{id} - u_{jd} + (n-1)x_{ijd} \le n-2
\qquad \forall d \in K,\; \forall i,j \in C,\; i \neq j.
\]

Daily route-length limit:
\[
\sum_{\substack{i \in L \\ }} \sum_{\substack{j \in L \\ j \neq i}} D_{ij}\, x_{ijd} \le M
\qquad \forall d \in K.
\]

\section*{Notes on Implementation}

\begin{itemize}
    \item The code uses zero-based Python indices for locations, where Python node \texttt{0} corresponds to business location $1$ (the depot). The mathematical formulation above uses the business numbering directly.
    \item The model is implemented as a mixed-integer linear program and solved with Gurobi through PuLP.
    \item The variables $u_{id}$ are created as continuous, not integer, with bounds $[1,n-1]$; this matches the implementation exactly.
    \item No $u$ variable is defined for the depot, and the MTZ constraints are imposed only for customer-customer arcs on each day.
    \item The formulation permits one of the two days to be unused: if $z_d=0$, then depot flow is zero and, by $y_{id}\le z_d$, no customer can be assigned to that day.
    \item Because every customer must be assigned exactly once across the two days, at least one day must be active.
\end{itemize}
